% ----------------------
\subsection{Section 84 (22--23 September 1832)}
% ----------------------
	\label{DC84}
	
	% -----
	\subsubsection{Historical Background}
	% -----
		\label{DC84-HB}
		
		% --
		\paragraph{JSP}
		% --
		
			``After JS relocated from Hiram to Kirtland, Ohio, on 12 September 1832, elders who had returned from preaching in the eastern United States came to JS to report on their proselytizing. While JS and these elders were `together in these seasons of joy,' a later history recounts, JS `inquired of the Lord and received' a revelation dated 22 and 23 September 1832. The beginning of the revelation identified its audience as JS and six elders, but partway through the revelation, the audience shifted to `Eleven high Priests save one.' Because JS was living in Newel K. Whitney's white store, the revelation was probably dictated in either the store's upstairs `translating room' or the upstairs `council room,' places where JS frequently worked.

			``The revelation was dictated over the course of two days, most likely beginning the evening of 22 September and continuing into the early morning hours of 23 September. Textual evidence indicates a pause in the dictation at some point on 23 September. The three existing manuscript copies of the revelation (one in the handwriting of Frederick G. Williams, another by Williams and JS, and one by John Whitmer) all contain a clear break between the phrase `for he is full of mercy Justice grace and truth and peace for ever and ever Amen' and the phrase `And again verily verily I say unto you it is expedient . . . ,' marking an interruption in the dictation. Whitmer's copy even adds `Received on the 23 day of September 1832' between those two lines. However, the three manuscripts also include `viz 23\supers{d.} day of September AD 1832' as a notation several pages before this break, indicating that material presented before the interruption was also dictated on 23 September. It may be that the dictation went into the early morning hours of 23 September, halted for a period of time, and then recommenced later that day.

			``The index of Revelation Book 2, one of the volumes in which this revelation was recorded, designated the revelation as one `explaining the two priest hoods and commissioning the Apostles to preach the gospel.' An understanding of priesthood was still developing among followers of JS, especially in terms of its connection to different offices in the church. The Book of Mormon indicated that authority from God was necessary to perform certain ordinances, such as baptism and conferring the gift of the Holy Ghost, which led JS and Oliver Cowdery to petition God for such authority as they worked together in translating that record. Later accounts indicate that in May 1829 and sometime thereafter, they received angelic visitations that provided them first with the authority to baptize and later with the authority to officiate in other ordinances. However, extant records up to June 1831 did not call such authority `priesthood'; that term---while appearing in both the Book of Mormon and in JS's Bible revision---did not appear in any other contemporary documents until the minutes of a June 1831 conference, which noted that several individuals `were ordained to the high Priesthood.' Moreover, the `Articles and Covenants' of the church explained the different duties of apostles, elders, priests, teachers, and deacons but did not explicitly associate these offices with the priesthood.

			``By late 1831, the high priesthood was understood to refer to both the office of high priest and to a broader authority. The office, according to an 11 November 1831 revelation, was superior to other offices in the church, just as the authority seemed to be the highest authority. This revelation stated that after the offices of deacon, teacher, priest, and elder came `the high Priest hood which is the greatest of all.' A history JS began writing around summer 1832 suggests that he had received two separate powers with different responsibilities. In that history, JS noted that `the ministring of---Aangels' gave him an authority that allowed him `to adminster the letter of the Gospel.' He also recorded receiving `the high Priesthood after the holy order of the son of the living God,' which gave him `power and ordinence from on high to preach the Gospel in the administration and demonstration of the spirit.'

			``This 22--23 September revelation similarly delineated the existence of two priesthoods: a greater priesthood that contains keys to the mysteries of the kingdom and to the knowledge of God, and a lesser priesthood holding the keys of the ministering of angels and of the gospel of repentance and baptism.11 The revelation traced the lineage of the two priesthoods, noting that the greater priesthood was held by Moses, who received it from a line of individuals (including Melchizedek) who had ultimately received it from God. Aaron, meanwhile, held the lesser priesthood, which was passed to his descendants until it reached John the Baptist. Both priesthoods, the revelation posited, are eternal, and men in the Church of Christ could become the `sons of Moses' and the `sons of Aaron' by receiving these priesthoods and `magnifying there calling.' Doing so would enable these `sons' to enter the temple that the Saints would construct in Missouri, allowing them to receive God's glory that would fill the temple. Some of this information seems a culmination of ideas first expressed in JS's revisions to passages in Hebrews, Genesis, and Exodus, which were completed in the six months leading up to September. The revelation also explained how the different offices in the church are connected to the two priesthoods. The offices of elder and bishop, it states, are `appendages' to the high priesthood, or to the office of high priest; the offices of teacher and deacon, meanwhile, are appendages to the lesser priesthood, or to the office of priest. High priests, elders, and priests, the revelation continued, have an obligation to travel to proclaim the gospel, while teachers and deacons are responsible for watching over the church in local congregations.

			``After providing a detailed explanation of the greater and lesser priesthoods, their offices, and their duties, the revelation emphasized the necessity of preaching to the world and provided an extensive discussion of many aspects of missionary work. Revelations from 1830, 1831, and 1832 called specific individuals on missions, but few revelations gave general procedural instructions about missionary work. This 22--23 September revelation, however, provided direction on who should serve missions, how they should serve, how they should receive sustenance while serving, what they should proclaim, and what would happen to those who did not accept their message. These instructions parallel New Testament accounts of the resurrected Jesus Christ's directions to the eleven apostles before his ascension into heaven. As in those accounts, Christ tells the high priests in this revelation that they have a responsibility to preach to all nations and to baptize those who believe. Signs will follow the believers, the revelation continues, and the Lord will go before them, just as Christ promised the ancient apostles that he would be with them always, `even unto the end of the world.' The revelation specifically refers to `Eleven high Priests save one' (perhaps a reference to the eleven apostles to whom Christ spoke in the New Testament) and calls these high priests `apostles' and `friends' of Jesus Christ. Although many had preached the gospel before this time, this revelation seemed to launch a more urgent and comprehensive missionary campaign, even including in its preaching assignments individuals such as Bishop Newel K. Whitney, who generally oversaw temporal, not spiritual, concerns.

			``Frederick G. Williams, who was serving as JS's scribe, probably recorded the revelation as JS dictated it. Williams also inscribed a copy of this revelation that was given to Whitney, stating that he was transcribing it `for N K, Whitney and Joseph the Seer.' Whitney and JS traveled together in New York and other states in October 1832, and the copy was likely made for the two men to carry with them on that trip. Williams and JS also made a copy of the revelation in Revelation Book 2, probably soon after the revelation was dictated. Parts of the revelation---including a `new song' on snake person themes, a condemnation of the church for neglecting the Book of Mormon, and the explanation of appendages to the greater and lesser priesthoods---were discussed in early 1833 issues of the church's periodical \textit{The Evening and the Morning Star}. A conference of high priests also wrote a letter in January 1833 calling church members in Missouri to repentance in conformance with instructions given in the revelation. Since at least six elders and ten high priests heard portions of the revelation while it was dictated, it is probable that information in it was also spread through word of mouth.''
			
	% -----
	\subsubsection{Versions}
	% -----
		\label{DC84-V}
		
		See the \href{https://drive.google.com/file/d/1goAvNz8jS4R8slYfMG_db55Ty2z_vmgK/view?usp=sharing}{sections comparison} document.
	
		\begin{compactdesc}
			\item[JSP] \hyperref[EMS-1832-JS-22-23-SEP-1832]{22--23 September 1832}
			\item[BC] ---
			\item[DC 1835] Section 4, 89--95
			\item[MS] 1:4 (August 1840), 77--81
			\item[TS] 5:19 (1 October 1844), 657--660
			\item[DC 1844] Section 4, 111--123
			\item[MS] 14:15 (5 June 1852), 228--231
		\end{compactdesc}

	% -----
	\subsubsection{Section 84:1} % *DC84-1 
	% -----
		\label{DC84-1}

		A REVELATION of Jesus Christ unto his servant Joseph Smith, Jun., and six elders, as they united their hearts and lifted their voices on high.

		% \begin{framed}
		% \scriptsize
			% \begin{compactdesc}
				% \item[JSP] 
				% % \item[BC] 
				% % \item[]
			% \end{compactdesc}
		% \end{framed}

	% -----
	\subsubsection{Section 84:2} % *DC84-2 
	% -----
		\label{DC84-2}

		Yea, the word of the Lord concerning his church, established in the last days for the restoration of his people, as he has spoken by the mouth of his prophets, and for the gathering of his saints to stand upon Mount Zion, which shall be the
			\footnote{%
				% \label{}%
				See \hyperref[ETHER-13-2]{Ether 13:2--12},
				}%
		city of New Jerusalem.

		% \begin{framed}
		% \scriptsize
			% \begin{compactdesc}
				% \item[JSP] 
				% % \item[BC] 
				% % \item[]
			% \end{compactdesc}
		% \end{framed}

	% -----
	\subsubsection{Section 84:3} % *DC84-3 
	% -----
		\label{DC84-3}

		Which city shall be built, beginning at the temple lot, which is appointed by the finger of the Lord, in the western boundaries of the State of Missouri, and dedicated by the hand of Joseph Smith, Jun., and others with whom the Lord was well pleased.

		% \begin{framed}
		% \scriptsize
			% \begin{compactdesc}
				% \item[JSP] 
				% % \item[BC] 
				% % \item[]
			% \end{compactdesc}
		% \end{framed}

	% -----
	\subsubsection{Section 84:4} % *DC84-4 
	% -----
		\label{DC84-4}

		Verily this is the word of the Lord, that the city New Jerusalem shall be built by the gathering of the saints, beginning at this place, even the place of the temple, which temple shall be reared in this generation.%
			\footnote{%
				% \label{}%
				I don't think this temple was ever built. I've been to see the temple lot area, what I remember is that it was just a field with a stone marking the position described here.
				}

		% \begin{framed}
		% \scriptsize
			% \begin{compactdesc}
				% \item[JSP] 
				% % \item[BC] 
				% % \item[]
			% \end{compactdesc}
		% \end{framed}

	% -----
	\subsubsection{Section 84:5} % *DC84-5 
	% -----
		\label{DC84-5}

		For verily this generation shall not all pass away until an house shall be built unto the Lord, and a cloud shall rest upon it, which cloud shall be even the glory of the Lord, which shall fill the house.

		\begin{framed}
		\scriptsize
			\begin{compactdesc}
				\item[JSP] for verely this generation shall not all pass away untill an house shalt be built unto the Lord and a cloud shall rest upon it which cloud shall be even the glory of the Lord which shall fill the house,
				% \item[BC] 
				% \item[]
			\end{compactdesc}
		\end{framed}

	% -----
	\subsubsection{Section 84:6} % *DC84-6 
	% -----
		\label{DC84-6}

		\footnote{%
			% \label{}%
			The transition from the previous verse to this one is incredibly jarring. What you have to realize is that everything after ``the sons of Moses'' in this verse, all the way to the end of \hyperref[DC84-30]{verse 30}, is like a parenthetical remark explaining different aspects of the priesthood. So if you want to read this without the parenthetical remark, it would start here and then continue in verse 30: ``And the sons of Moses...and also the sons of Aaron shall offer an acceptable offering and sacrifice in the house of the Lord...''
			}%
		And the sons of Moses,%
			\footnote{%
				% \label{}%
				The ``sons of Moses'' must be people who have received the Melchizedek priesthood, or at least something more than the priesthood of Aaron, as is made clear in \hyperref[DC84-17]{verses 17--18}, where the revelation makes a distinction between them.
				}
		according to the Holy Priesthood%
			\footnote{%
				% \label{}%
				I want to read the passage this way: ``And those who are the sons of Moses \textit{according to the Holy Priesthood}.'' In other words, ``according to the Holy Priesthood'' explains the respect in which these individuals are sons of Moses. This goes along with \hyperref[DC84-32]{verse 32}, which says that those listening are these sons.	
				}
		which he received under the hand of his father-in-law, \hyperref[JETHRO]{Jethro};%
			\footnote{%
				% \label{}%
				There follows here a lineage of the priesthood being described; to me the whole thing sounds made up. I don't doubt it could have been given by revelation, and that there could have been people with these names, but this seems like the kind of thing that JS would make up.
				}

		\begin{framed}
		\scriptsize
			\begin{compactdesc}
				\item[JSP] and the sons of Moses according to the holy Priesthood which he received under the father in Law Jethro,
				% \item[BC] 
				% \item[]
			\end{compactdesc}
		\end{framed}

	% -----
	\subsubsection{Section 84:7} % *DC84-7 
	% -----
		\label{DC84-7}

		And \hyperref[JETHRO]{Jethro} received it under the hand of \hyperref[CALEB]{Caleb};%
			\footnote{%
				% \label{}%
				There are two ``Caleb''s listed in NIEBC. The first appears in Numbers and Joshua, and the second apperas in 1 Chronicles. If this reference here is meant to speak of one of these Calebs, it must be referring to the first.
				}

		% \begin{framed}
		% \scriptsize
			% \begin{compactdesc}
				% \item[JSP] 
				% % \item[BC] 
				% % \item[]
			% \end{compactdesc}
		% \end{framed}

	% -----
	\subsubsection{Section 84:8} % *DC84-8 
	% -----
		\label{DC84-8}

		And \hyperref[CALEB]{Caleb} received it under the hand of Elihu;%
			\footnote{%
				% \label{}%
				NIEBC lists five people named ``Elihu.'' None appear early enough, I think, to have been the one spoken of here, who gave the priesthood to Caleb and received it from Jeremy.
				}

		% \begin{framed}
		% \scriptsize
			% \begin{compactdesc}
				% \item[JSP] 
				% % \item[BC] 
				% % \item[]
			% \end{compactdesc}
		% \end{framed}

	% -----
	\subsubsection{Section 84:9} % *DC84-9 
	% -----
		\label{DC84-9}

		And Elihu under the hand of Jeremy;%
			\footnote{%
				% \label{}%
				NIEBC lists nine people in the OT with the name ``Jeremiah'' (there are none with the name ``Jeremy''). The first of these is the prophet Jeremiah, and he's by far the best known; there are eight others, though none of them, including the big Jeremiah, appear early enough to be the person referred to here.
				}

		% \begin{framed}
		% \scriptsize
			% \begin{compactdesc}
				% \item[JSP] 
				% % \item[BC] 
				% % \item[]
			% \end{compactdesc}
		% \end{framed}

	% -----
	\subsubsection{Section 84:10} % *DC84-10 
	% -----
		\label{DC84-10}

		And Jeremy under the hand of Gad;%
			\footnote{%
				% \label{}%
				Jacob's seventh son, born to Zilpah, was named Gad; I don't see how that could be the person spoken of here, since the timing doesn't seem to work. But maybe.
				}

		% \begin{framed}
		% \scriptsize
			% \begin{compactdesc}
				% \item[JSP] 
				% % \item[BC] 
				% % \item[]
			% \end{compactdesc}
		% \end{framed}

	% -----
	\subsubsection{Section 84:11} % *DC84-11 
	% -----
		\label{DC84-11}

		And Gad under the hand of Esaias;%
			\footnote{%
				% \label{}%
				This is a term for ``Isaiah,'' apparently.
				}

		% \begin{framed}
		% \scriptsize
			% \begin{compactdesc}
				% \item[JSP] 
				% % \item[BC] 
				% % \item[]
			% \end{compactdesc}
		% \end{framed}

	% -----
	\subsubsection{Section 84:12} % *DC84-12 
	% -----
		\label{DC84-12}

		And Esaias received it under the hand of God.

		% \begin{framed}
		% \scriptsize
			% \begin{compactdesc}
				% \item[JSP] 
				% % \item[BC] 
				% % \item[]
			% \end{compactdesc}
		% \end{framed}

	% -----
	\subsubsection{Section 84:13} % *DC84-13 
	% -----
		\label{DC84-13}

		Esaias also lived in the days of Abraham, and was blessed of him---

		% \begin{framed}
		% \scriptsize
			% \begin{compactdesc}
				% \item[JSP] 
				% % \item[BC] 
				% % \item[]
			% \end{compactdesc}
		% \end{framed}

	% -----
	\subsubsection{Section 84:14} % *DC84-14 
	% -----
		\label{DC84-14}

		Which Abraham received the priesthood from Melchizedek,%
			\footnote{%
				% \label{}%
				See \hyperref[JST-OTR1-GENESIS-14]{OTR1 of Genesis 14} for a lot more here about Melchizedek.
				}
		who received it through the lineage of his fathers,%
			\footnote{%
				% \label{}%
				See this ``lineage'' language also in \hyperref[DC86-8]{DC 86:8}.
				}
		even till Noah;

		% \begin{framed}
		% \scriptsize
			% \begin{compactdesc}
				% \item[JSP] 
				% % \item[BC] 
				% % \item[]
			% \end{compactdesc}
		% \end{framed}

	% -----
	\subsubsection{Section 84:15} % *DC84-15 
	% -----
		\label{DC84-15}

		And from Noah till Enoch, through the lineage of their fathers;

		% \begin{framed}
		% \scriptsize
			% \begin{compactdesc}
				% \item[JSP] 
				% % \item[BC] 
				% % \item[]
			% \end{compactdesc}
		% \end{framed}

	% -----
	\subsubsection{Section 84:16} % *DC84-16 
	% -----
		\label{DC84-16}

		And from Enoch to Abel, who was slain by the conspiracy of his brother, who received the priesthood by the commandments of God, by the hand of his father Adam, who was the first man---

		% \begin{framed}
		% \scriptsize
			% \begin{compactdesc}
				% \item[JSP] 
				% % \item[BC] 
				% % \item[]
			% \end{compactdesc}
		% \end{framed}

	% -----
	\subsubsection{Section 84:17} % *DC84-17 
	% -----
		\label{DC84-17}

		Which priesthood continueth in the church of God in all generations, and is without beginning of days or end of years.%
			\footnote{%
				% \label{}%
				For this ``beginning...years'' phrase at the end, see \hyperref[END-PHRASES]{here}.
				}

		% \begin{framed}
		% \scriptsize
			% \begin{compactdesc}
				% \item[JSP] 
				% % \item[BC] 
				% % \item[]
			% \end{compactdesc}
		% \end{framed}

	% -----
	\subsubsection{Section 84:18} % *DC84-18 
	% -----
		\label{DC84-18}

		And the Lord confirmed a priesthood also upon Aaron and his seed, throughout all their generations, which priesthood also continueth and abideth forever with the priesthood which is after the holiest \hyperref[ORDER]{order}%
			\footnote{%
				% \label{}%
				This was hard for me to believe, but this ``order'' here is the only occurrence of this word in this section. It occurs some 14 times in \hyperref[DC107]{section 107}.
				}
		of God.%
			\footnote{%
				% \label{}%
				This last part of the verse is referring to the Melchizedek priesthood, and calls it the one ``after the holiest order of God.'' See \hyperref[DC107-3]{DC 107:3}.
				}

		\begin{framed}
		\scriptsize
			\begin{compactdesc}
				\item[JSP] and the Lord confirmed a priesthood also upon Aaron and his seed throughout all the generations of the Jews. which priesthood also continueth and abideth for ever with the Priesthood which is after the holiest order of God, 
				% \item[BC] 
				% \item[]
			\end{compactdesc}
		\end{framed}

	% -----
	\subsubsection{Section 84:19} % *DC84-19 
	% -----
		\label{DC84-19}

		And this greater priesthood administereth the gospel and holdeth the \hyperref[KEY]{key} of the \hyperref[MYSTERY]{mysteries} of the kingdom, even the \hyperref[KEY]{key} of the knowledge%
			\footnote{%
				% \label{}%
				Subjective or objective genitive? Maybe both? Subjective: ``the key which acts as God's knowledge; the key which God's knowledge is.'' Objective: ``the key to God's knowledge; the key to knowledge about God.''
				}
		of God.%
			\footnote{%
				% \label{}%
				See \hyperref[DC107-18]{DC 107:18--19}.
				}

		\begin{framed}
		\scriptsize
			\begin{compactdesc}
				\item[JSP] and this greater Priesthood adminestereth the gospel and holdeth the key of the misteries of the kingdom, even the key of the knowledge of God
				% \item[BC] 
				% \item[]
			\end{compactdesc}
		\end{framed}

	% -----
	\subsubsection{Section 84:20} % *DC84-20 
	% -----
		\label{DC84-20}

			\footnote{%
				% \label{}%
				For this verse and the several after it, compare the very important verse \hyperref[2TIMOTHY-3-5]{2 Timothy 3:5}. The ``power'' spoken of these is \hyperref[DUNAMIS]{\grl{δύναμις}}, not \hyperref[EXOUSIA]{\grl{ἐξουσία}}. These verses maintain the distinction in speaking of ``authority'' in the next verse---the ``power'' here is \gr{δύναμις}.
				}%
		Therefore,
		\footnote{%
			% \label{}%
			I've also been very impressed by the relationship between these few verses here and \hyperref[ROMAN-1-16]{Romans 1:16--20}, so much so that it almost seems like these verses are a commentary on parts of those.
			}%
		in the \hyperref[ORDINANCE]{ordinances} thereof,%
			\footnote{%
				% \label{}%
				Of the Melchizedek priesthood.
				}
		the power of \hyperref[GODLINESS]{godliness}%
			\footnote{%
				% \label{}%
				The text does NOT say that the ``power of \textit{God}'' is manifest; it says that the power of \textit{godliness} is manifest. It is the ordinances of the priesthood---including things like resurrection, I'm thinking---that the power of godliness or godly living or a spiritual life is manifest. Without those ordinances, or absent the real power of the priesthood, you won't be able to discover the true power of living that way.
				}
		is \hyperref[MANIFEST]{manifest}.%
			\footnote{%
				% \label{}%
				The word here is not ``manifested'' but ``manifest.'' The difference matters because ``manifest'' is an adjective while the past participle is ``manifested.'' So the question here is, is this a passive verb construction, ``the power of godliness is manifest\textit{ed},'' or is it a predicate adjective, ``the power of godliness is manifest [adjective]''? I believe it's the latter. ``Manifested'' appears five times in the DC, including in three passive constructions (\hyperref[DC20-5]{20:5}, \hyperref[DC90-14]{90:14}, \hyperref[DC115-18]{115:18}), and so at least some passages respect the distinction. The OED for ``\hyperref[MANIFEST-OED]{manifest},'' in meaning \#2.b, seems to be the primary meaning/use we would see here for a passive verb, but all the examples listed (of passive voice) use ``manifested.'' That isn't sure proof, but I think the evidence is heavily on the side of ``manifest'' in this verse being an adjective and not part of a passive verb construction. I will say, though, that the use in the next verse looks more like a passive verb, though I still want to read it as an adjective. Compare \hyperref[DC93-31]{DC 93:31}.
				}

		\begin{framed}
		\scriptsize
			\begin{compactdesc}
				\item[JSP] therefore in the ordinences thereof the power of Godliness is manifest
				% \item[BC] 
				% \item[]
			\end{compactdesc}
		\end{framed}

	% -----
	\subsubsection{Section 84:21} % *DC84-21 
	% -----
		\label{DC84-21}

		And without the ordinances thereof, and the authority of the priesthood, the power of \hyperref[GODLINESS]{godliness} is not manifest unto men in the flesh;%
			\footnote{%
				% \label{}%
				Compare \hyperref[DC88-34]{DC 88:34}.
				}

		\begin{framed}
		\scriptsize
			\begin{compactdesc}
				\item[JSP] and without the ordinences thereof, and the authority of the Priesthood, the power of Godliness is not manifest unto man in the flesh,
				% \item[BC] 
				% \item[]
			\end{compactdesc}
		\end{framed}

	% -----
	\subsubsection{Section 84:22} % *DC84-22 
	% -----
		\label{DC84-22}

		For without this%
			\footnote{%
				% \label{}%
				Antecedent of this pronoun? Could be the ``ordinances,'' the ``authority'' of the priesthood, the ``power'' of godliness, or a combination of all three: ``For without this [system], no man can see the face of God.''
				}
		no man can see the face of God, even the Father, and live.

		\begin{framed}
		\scriptsize
			\begin{compactdesc}
				\item[JSP] for without this no man can see the face of God even the father and live,
				% \item[BC] 
				% \item[]
			\end{compactdesc}
		\end{framed}

	% -----
	\subsubsection{Section 84:23} % *DC84-23 
	% -----
		\label{DC84-23}

			\footnote{%
				% \label{}%
				For this verse and the next several, see \hyperref[EXODUS-34-1]{Exodus 34:1--2}, especially the OTR2.
				}%
		Now this Moses plainly taught to the children of Israel in the wilderness, and sought diligently to sanctify his people that they might behold the face of God;

		\begin{framed}
		\scriptsize
			\begin{compactdesc}
				\item[JSP] now this Moses plainly taught to the children of Israel in the wilderness, and saught diligently to sanctify his people that they might behold the face of God,
				% \item[BC] 
				% \item[]
			\end{compactdesc}
		\end{framed}

	% -----
	\subsubsection{Section 84:24} % *DC84-24 
	% -----
		\label{DC84-24}

		But they hardened their hearts and could not endure his presence; therefore, the Lord in his wrath, for his anger was kindled against them, swore that they should not enter into his rest while in the wilderness, which rest%
			\footnote{%
				% \label{}%
				Compare \hyperref[MATTHEW-11-28]{Matthew 11:28--30}, \hyperref[DC51-19]{DC 51:19}, \hyperref[DC54-10]{DC 54:10}, \hyperref[DC124-86]{DC 124:86}, 
				}
		is the fulness of his glory.

		\begin{framed}
		\scriptsize
			\begin{compactdesc}
				\item[JSP] but they hardened ther hearts and could not endure his presence therefore the Lord in his wrath (for his anger was kindled against them) swore that they should not enter into his rest, which rest is the fulness of his glory while in the wilderness,
				% \item[BC] 
				% \item[]
			\end{compactdesc}
		\end{framed}

	% -----
	\subsubsection{Section 84:25} % *DC84-25 
	% -----
		\label{DC84-25}

		Therefore, he took Moses out of their midst, and the Holy Priesthood also;

		\begin{framed}
		\scriptsize
			\begin{compactdesc}
				\item[JSP] therefore he took Moses out of there midst and the holy Priesthood also,
				% \item[BC] 
				% \item[]
			\end{compactdesc}
		\end{framed}

	% -----
	\subsubsection{Section 84:26} % *DC84-26 
	% -----
		\label{DC84-26}

		And the lesser priesthood continued, which priesthood holdeth the key of the ministering of angels and the preparatory gospel;%
			\footnote{%
				% \label{}%
				Compare the description of the Aaronic priesthood in \hyperref[DC13-1]{DC 13}, and also \hyperref[DC107-20]{107:20}.
				}

		\begin{framed}
		\scriptsize
			\begin{compactdesc}
				\item[JSP] and the lesser Priesthood continued, which Priesthood holdeth the keys of the ministring of Angels and the preparitory gospel,
				% \item[BC] 
				% \item[]
			\end{compactdesc}
		\end{framed}

	% -----
	\subsubsection{Section 84:27} % *DC84-27 
	% -----
		\label{DC84-27}

		Which gospel is the gospel of repentance and of baptism, and the remission of sins, and the law of carnal commandments, which the Lord in his wrath caused to continue with the house of Aaron among the children of Israel until John, whom God raised up, being filled with the Holy Ghost from his mother's womb.

		\begin{framed}
		\scriptsize
			\begin{compactdesc}
				\item[JSP] which gospel is the gospel of repentence and of Baptism, and the remission of sins, and the Law of carnal commandments--- which the lord in his wrath caused to continue with the house of Aaron among the children of Israel until John whom God raised up being fillid with the holy ghost from his Mothers womb,
				% \item[BC] 
				% \item[]
			\end{compactdesc}
		\end{framed}

	% -----
	\subsubsection{Section 84:28} % *DC84-28 
	% -----
		\label{DC84-28}

		For he was baptized while he was yet in his childhood, and was ordained by the angel of God at the time he was eight days old unto this power, to overthrow the kingdom of the Jews, and to make straight the way of the Lord before the face of his people, to prepare them for the coming of the Lord, in whose hand is given all power.

		\begin{framed}
		\scriptsize
			\begin{compactdesc}
				\item[JSP] for he was baptised while he was yet in \sout{his} the \sout{mothers} womb and was ordained by the Angel of God at the time he was eight days old unto this power to overthrow the kingdom of the Jews and to make straight the way of the Lord before the face of his people to prepare them for the coming of the Lord in whose hand is given all power,
				% \item[BC] 
				% \item[]
			\end{compactdesc}
		\end{framed}

	% -----
	\subsubsection{Section 84:29} % *DC84-29 
	% -----
		\label{DC84-29}

		And again, the offices of elder and bishop are necessary appendages belonging unto the high priesthood.

		\begin{framed}
		\scriptsize
			\begin{compactdesc}
				\item[JSP] and again, the offices of Elder \& Bishop are necessary appendages belon[g]ing unto the high Priesthood,
				% \item[BC] 
				% \item[]
			\end{compactdesc}
		\end{framed}

	% -----
	\subsubsection{Section 84:30} % *DC84-30 
	% -----
		\label{DC84-30}

		And again, the offices of teacher and deacon are necessary appendages belonging to the lesser priesthood, which priesthood was confirmed upon Aaron and his sons.

		\begin{framed}
		\scriptsize
			\begin{compactdesc}
				\item[JSP] and again the offices of Teacher and Deacon are necessary appendages belonging to the lesser Priesthood, which priesthood was confirmed upon Aaron and his sons
				% \item[BC] 
				% \item[]
			\end{compactdesc}
		\end{framed}

	% -----
	\subsubsection{Section 84:31} % *DC84-31 
	% -----
		\label{DC84-31}

			\footnote{%
				% \label{}%
				With the end of the previous verse we get the end of the long parenthetical remark that interrupts the discussion in \hyperref[DC84-6]{verse 6}. So the right way to see what's going on here is that verse 31 is picking up on what was happening at the start of verse 6, and so it's clearly related to the temple.
				}%
		Therefore, as I said concerning the sons of Moses---for the sons of Moses and also the sons of Aaron shall offer an acceptable offering and sacrifice in the house of the Lord, which house shall be built unto the Lord in this generation, upon the consecrated spot as I have appointed---

		\begin{framed}
		\scriptsize
			\begin{compactdesc}
				\item[JSP] therefore as I said, concerning the Sons of Moses, for the sons of Moses, and also the sons of Aaron shall offer an aaceptable offering and sacrifice in the house of the Lord which house \sout{shall} shalt be built unto the Lord in this generation upon the consecrated spot as I have appointed
				% \item[BC] 
				% \item[]
			\end{compactdesc}
		\end{framed}

	% -----
	\subsubsection{Section 84:32} % *DC84-32 
	% -----
		\label{DC84-32}

		And the sons of Moses and of Aaron shall be filled with the glory of the Lord, upon Mount Zion in the Lord's house, whose sons are ye; and also many whom I have called and sent forth to build up my church.

		\begin{framed}
		\scriptsize
			\begin{compactdesc}
				\item[JSP] and the sons of Moses, and of Aaron shall be filled with the glory of the Lord upon mount Zion in the Lords house whose sons are ye, and also many whom I have called and sent forth to build up my church
				% \item[BC] 
				% \item[]
			\end{compactdesc}
		\end{framed}

	% -----
	\subsubsection{Section 84:33} % *DC84-33 
	% -----
		\label{DC84-33}

			\footnote{%
				% \label{}%
				This verse begins an explanatory comment that elaborates on how someone becomes a son of Moses and Aaron, and what that means. Note, though, that it also introduces the idea of being the ``seed of Abraham,'' which is new, and has not appeared (I think) so far in this revelation.
				}%
		For \hyperref[WHOSO]{whoso} is faithful unto the obtaining these two priesthoods%
			\footnote{%
				% \label{}%
				Must be referring to the Melchizedek priesthood and Aaronic priesthood.
				}
		of which I have spoken, and the \hyperref[MAGNIFY]{magnifying}%
			\footnote{%
				% \label{}%
				See \hyperref[ROMANS-11-13]{Romans 11:13}, 
				}
		their calling,%
			\footnote{%
				% \label{}%
				See \hyperref[DC24-9]{DC 24:9}.
				}
		are sanctified by the Spirit unto the renewing%
			\footnote{%
				% \label{}%
				See also \hyperref[ROMANS-12-2]{Romans 12:2}, \hyperref[2CORINTHIANS-4-10]{2 Corinthians 4:10--12}, \hyperref[2CORINTHIANS-4-16]{2 Corinthians 4:16}, \hyperref[EPHESIANS4-23]{Ephesians 4:23}, \hyperref[DC84-80]{DC 84:80},
				}
		of their
			\footnote{%
				% \label{}%
				This is the only place in the scriptures where you get language about ``renewing bodies.'' For a similar idea of receiving strength or renewal in service, see \hyperref[MOSIAH-2-11]{Mosiah 2:11}, \hyperref[DC24-7]{DC 24:7, 12},
				}%
		bodies.

		\begin{framed}
		\scriptsize
			\begin{compactdesc}
				\item[JSP] for whoso is faithful unto the attaining of these two Priesthoods of which I have spoken and the magnifying there calling are sanctified by the spir[i]t unto the renewing of there bodies
				% \item[BC] 
				% \item[]
			\end{compactdesc}
		\end{framed}

	% -----
	\subsubsection{Section 84:34} % *DC84-34 
	% -----
		\label{DC84-34}

			\footnote{%
				% \label{}%
				Very important omission in this version, compared to the JSP---it's missing a ``that,'' which connects this verse as a purpose or result clause following the previous one.
				}%
		They become the sons of Moses and of Aaron and the seed of Abraham,%
			\footnote{%
				% \label{}%
				I don't think there's anything so far in this revelation which refers even indirectly to a relation with Abraham; this seems to be something new, though note that Abraham does appear in this section, in \hyperref[DC84-13]{verses 13--14}. See also \hyperref[DC132-29]{DC 132:29--33} and \hyperref[DC132-49]{49--50}, but especially \hyperref[ABRAHAM-2-8]{Abraham 2:8--11}. The ``seed'' is the key idea.
				}
		and the church and kingdom, and the \hyperref[ELECTION]{elect} of God.

		\begin{framed}
		\scriptsize
			\begin{compactdesc}
				\item[JSP] that they become the sons of Moses and of Aaron and the seed of Abraham and the church and kingdom and the elect of God
				% \item[BC] 
				% \item[]
			\end{compactdesc}
		\end{framed}

	% -----
	\subsubsection{Section 84:35} % *DC84-35 
	% -----
		\label{DC84-35}

		And also all they who receive this priesthood%
			\footnote{%
				% \label{}%
				\hyperref[DC84-33]{Verse 33} above mentions two priesthoods, which seem to be the Melchizedek priesthood and the Aaronic priesthood. But now we're talking about ``this priesthood,'' singular. President Openshaw's explanation for this was that the text is now speaking about a different order of priesthood, perhaps called ``patriarchal,'' which gets introduced indirectly with the reference to the ``seed of Abraham'' in the previous verse. I think this may be right, though I'm not sure about the ``patriarchal'' title; it would be enough to think of it as a distinct order. Note that, in \hyperref[DC84-13]{verses 13--14}, we hear about Abraham receiving ``the priesthood'' from Melchizedek, who in turn received it ``through the lineage of his fathers'' (note that that last bit supports the ``patriarchal'' idea). So here's one way to try to organize all of this---there are three orders of the priesthood being discussed here, with three individuals or sets of individuals as representatives of them: Aaron representing the Aaronic order, Moses representing the Melchizedek order, and Abraham and maybe Melchizedek himself representing the higher order now being introduced, possibly ``patriarchal.'' All of this, however, needs to be understood in the light of JS's later comments, including one from \hyperref[EMS-1841-JS-5-JAN-1841-WC]{5 January 1841}, where he says that ``All priesthood is Melchizedeck; but there are different portions or degrees of it.'' We can't forget that---it's crucial to understanding the priesthood.
				}
		receive me, saith the Lord;

		\begin{framed}
		\scriptsize
			\begin{compactdesc}
				\item[JSP] and also all they who receive this Priesthood receiveth me saith the Lord
				% \item[BC] 
				% \item[]
			\end{compactdesc}
		\end{framed}

	% -----
	\subsubsection{Section 84:36} % *DC84-36 
	% -----
		\label{DC84-36}

		For he that receiveth my servants receiveth me;%
			\footnote{%
				% \label{}%
				See also \hyperref[MATTHEW-10-40]{Matthew 10:40--41}; \hyperref[DC84-89]{DC 84:89}; \hyperref[DC99-2]{DC 99:2--4}; 
				}

		\begin{framed}
		\scriptsize
			\begin{compactdesc}
				\item[JSP] for he that receiveth my servants receveth me,
				% \item[BC] 
				% \item[]
			\end{compactdesc}
		\end{framed}

	% -----
	\subsubsection{Section 84:37} % *DC84-37 
	% -----
		\label{DC84-37}

		And he that receiveth me receiveth my Father;

		\begin{framed}
		\scriptsize
			\begin{compactdesc}
				\item[JSP] and he that receiveth me receiveth my father,
				% \item[BC] 
				% \item[]
			\end{compactdesc}
		\end{framed}

	% -----
	\subsubsection{Section 84:38} % *DC84-38 
	% -----
		\label{DC84-38}

		And he that receiveth my Father receiveth my Father's kingdom; therefore all that my Father hath shall be given unto him.

		\begin{framed}
		\scriptsize
			\begin{compactdesc}
				\item[JSP] and he that receiveth my father, receiveth my fathers kingdom, therefore all that my father hath shall be given unto him
				% \item[BC] 
				% \item[]
			\end{compactdesc}
		\end{framed}

	% -----
	\subsubsection{Section 84:39} % *DC84-39 
	% -----
		\label{DC84-39}

		And this is according to the oath and covenant which belongeth to the priesthood.

		\begin{framed}
		\scriptsize
			\begin{compactdesc}
				\item[JSP] and this is according to the oath and the covenant which belongeth to the Priesthood,
				% \item[BC] 
				% \item[]
			\end{compactdesc}
		\end{framed}

	% -----
	\subsubsection{Section 84:40} % *DC84-40 
	% -----
		\label{DC84-40}

			\footnote{%
				% \label{}%
				Roughly about here, though it's maybe not exact, is the place where I would say that the explanatory comment started in \hyperref[DC84-33]{verse 33} ends. Or I mean, it ends in the previous verse.
				}%
		Therefore, all those who receive the priesthood,%
			\footnote{%
				% \label{}%
				Notice that in this verse and the previous one we are now talking about ``\textit{the} priesthood,'' not ``these two priesthoods'' or ``this priesthood.'' We seem to have moved to talking about it in more general terms.
				}
		receive this oath and covenant of my Father, which he cannot break,%
			\footnote{%
				% \label{}%
				Compare \hyperref[DC82-10]{DC 82:10}.
				} 
		neither can it be moved.

		\begin{framed}
		\scriptsize
			\begin{compactdesc}
				\item[JSP] therefore all those. who receive the Priesthood, receiveth this oath and covenant of my father which he cannot break neither can it be mooved,
				% \item[BC] 
				% \item[]
			\end{compactdesc}
		\end{framed}

	% -----
	\subsubsection{Section 84:41} % *DC84-41 
	% -----
		\label{DC84-41}

		But whoso breaketh this covenant after he hath received it, and altogether turneth therefrom, shall not have forgiveness of sins in this world nor in the world to come.%
			\footnote{%
				% \label{}%
				Possibly a reference to the \hyperref[SONSOFPERDITION]{sons of perdition}? I also read this as related to \hyperref[ALMA-32-19]{Alma 32:19} idea of ``knowing the will of God concerning you,'' but I think what we're talking about in this verse is a greater level of accountability.
				}

		\begin{framed}
		\scriptsize
			\begin{compactdesc}
				\item[JSP] but whoso breaketh this covenant after he hath received it, and altogether turneth therefrom shall not have forgivness in this world nor in the world to come
				% \item[BC] 
				% \item[]
			\end{compactdesc}
		\end{framed}

	% -----
	\subsubsection{Section 84:42} % *DC84-42 
	% -----
		\label{DC84-42}

		And wo unto all those who come not unto this priesthood%
			\footnote{%
				% \label{}%
				Back to the ``this priesthood'' language, and specifically, the Lord is now talking about the one that the audience for this revelation has received.
				}
		which ye have received, which I now confirm%
			\footnote{%
				% \label{}%
				``Confirm'' reminds me of the \hyperref[TEMPLE-ENDOWMENT-INIT-PRE-2005-ANOINT]{anointing} in the initiatory ceremony.
				}
		upon you who are present this day, by mine own voice out of the heavens; and even I have given the heavenly hosts and mine angels charge concerning you.

		\begin{framed}
		\scriptsize
			\begin{compactdesc}
				\item[JSP] and all those that come not unto this Priesthood, which ye have received, which I now confirm upon you who are present this day viz the 23\supers{d.} day of September AD 1832 Eleven high Priests save one by by mine own voice out of the heavens and even I have given the heavenly hosts and mine Angels charge concerning you,
				% \item[BC] 
				% \item[]
			\end{compactdesc}
		\end{framed}

	% -----
	\subsubsection{Section 84:43} % *DC84-43 
	% -----
		\label{DC84-43}

		And I now give unto you a commandment to beware concerning yourselves, to give diligent heed to the words of eternal life.%
			\footnote{%
				% \label{}%
				For ``words of eternal life,'' see \hyperref[ETERNITY-PHRASES]{here}.
				}

		\begin{framed}
		\scriptsize
			\begin{compactdesc}
				\item[JSP] and I now give unto you a commandment to beware concerning yourselves to give heed dilligently to the words of eternal life
				% \item[BC] 
				% \item[]
			\end{compactdesc}
		\end{framed}

	% -----
	\subsubsection{Section 84:44} % *DC84-44 
	% -----
		\label{DC84-44}

			\footnote{%
				% \label{}%
				Notice that this verse begins with a ``for,'' marking it as an explanatory comment on what came before; in fact this is a really poor verse division and verses 43 and 44 should be one verse.
				}%
		For you shall live by every word that proceedeth forth from the mouth of God.%
			\footnote{%
				% \label{}%
				For me, this verse is maybe the premier example of living in the flow of the Spirit, and of making real the promise in the sacrament prayers that we will always have the Spirit to be with us. Notice that we're supposed to ``give heed'' to the words of eternal life, per the previous verse; Christ has these words (\hyperref[JOHN-6-68]{John 6:68}), and both angels and the Holy Ghost can deliver these words to us (\hyperref[2NEPHI-32-3]{2 Nephi 32:3--5}). See also \hyperref[1NEPHI-17-13]{1 Nephi 17:13}, \hyperref[DC6-14]{DC 6:14--15}; \hyperref[DC28-15]{DC 28:15}; \hyperref[DC97-1]{DC 97:1} (which allows you to connect the ``voice'' and ``spirit'' to these verses here); \hyperref[DC98-11]{DC 98:11}; 
				}

		\begin{framed}
		\scriptsize
			\begin{compactdesc}
				\item[JSP] for you shall live by evry word that procedeth forth from the mouth of God
				% \item[BC] 
				% \item[]
			\end{compactdesc}
		\end{framed}

	% -----
	\subsubsection{Section 84:45} % *DC84-45 
	% -----
		\label{DC84-45}

			\footnote{%
				% \label{}%
				This verse begins an explanatory comment on which elaborates on verses 43--44. This comment looks like it goes through \hyperref[DC84-48]{verse 48}.
				}%
		For%
			\footnote{%
				% \label{}%
				Note the connection between this verse and the next one and \hyperref[DC88-11]{DC 88:11--13}. There are also clear connections to \hyperref[DC93]{DC 93}.
				}
		the word of the Lord is truth, and whatsoever is truth is light, and whatsoever is light is Spirit, even the Spirit of Jesus Christ.

		\begin{framed}
		\scriptsize
			\begin{compactdesc}
				\item[JSP] for the word of the Lord is truth and whatsoever is truth is light, and whatsoever is light is spirit even the spirit of Jesus Christ,
				% \item[BC] 
				% \item[]
			\end{compactdesc}
		\end{framed}

	% -----
	\subsubsection{Section 84:46} % *DC84-46 
	% -----
		\label{DC84-46}

		And the Spirit giveth light to every man that cometh into the world; and the Spirit enlighteneth every man through the world, that hearkeneth to the voice of the Spirit.

		\begin{framed}
		\scriptsize
			\begin{compactdesc}
				\item[JSP] and the spirit giveth light to evry man that cometh into the world, and the spirit enlightneth evry man through the world that harkneth to the voice of the spirit,
				% \item[BC] 
				% \item[]
			\end{compactdesc}
		\end{framed}

	% -----
	\subsubsection{Section 84:47} % *DC84-47 
	% -----
		\label{DC84-47}

		And every one that hearkeneth to the voice of the Spirit cometh unto God, even the Father.

		\begin{framed}
		\scriptsize
			\begin{compactdesc}
				\item[JSP] and evry one that harkneth to the voice of the spirit cometh unto God even the father
				% \item[BC] 
				% \item[]
			\end{compactdesc}
		\end{framed}

	% -----
	\subsubsection{Section 84:48} % *DC84-48 
	% -----
		\label{DC84-48}

			\footnote{%
				% \label{}%
				This verse seems to be the last one in the explanatory comment to \hyperref[DC84-43]{verses 43--44} which begins in verse 45. Notice that we're talking again about ``confirming'' things, which we last did in verse 42.
				}%
		And the Father%
			\footnote{%
				% \label{}%
				I think it's extremely significant that the Father is doing the teaching here, though I'm not totally sure why yet. See \hyperref[JOHN-6-45]{John 6:45} and notes there for some important commentary---according to that verse we are to hear and learn from the Father; the Father also does the ``draw''-ing in \hyperref[JOHN-6-44]{verse 44}. See \hyperref[JOHN-8-28]{John 8:28}, \hyperref[MORMON-5-17]{Mormon 5:17}, \hyperref[ETHER-6-17]{Ether 6:17}, 
				}
		teacheth him of the covenant which he has renewed and confirmed upon you, which is confirmed upon you for your sakes, and not for your sakes only, but for the sake of the whole world.

		\begin{framed}
		\scriptsize
			\begin{compactdesc}
				\item[JSP] and the father teacheth him of the covenant which he hath renewed and confirmed upon you which is confirmed upon you for your sakes \sout{and} and not for yours only, but for the sake of the whole world,
				% \item[BC] 
				% \item[]
			\end{compactdesc}
		\end{framed}

	% -----
	\subsubsection{Section 84:49} % *DC84-49 
	% -----
		\label{DC84-49}

		And the whole world lieth in sin, and groaneth under darkness and under the bondage of sin.

		% \begin{framed}
		% \scriptsize
			% \begin{compactdesc}
				% \item[JSP] 
				% % \item[BC] 
				% % \item[]
			% \end{compactdesc}
		% \end{framed}

	% -----
	\subsubsection{Section 84:50} % *DC84-50 
	% -----
		\label{DC84-50}

		And by this you may know they are under the bondage of sin, because they come not unto me.

		% \begin{framed}
		% \scriptsize
			% \begin{compactdesc}
				% \item[JSP] 
				% % \item[BC] 
				% % \item[]
			% \end{compactdesc}
		% \end{framed}

	% -----
	\subsubsection{Section 84:51} % *DC84-51 
	% -----
		\label{DC84-51}

		For whoso cometh not unto me is under the bondage of sin.

		% \begin{framed}
		% \scriptsize
			% \begin{compactdesc}
				% \item[JSP] 
				% % \item[BC] 
				% % \item[]
			% \end{compactdesc}
		% \end{framed}

	% -----
	\subsubsection{Section 84:52} % *DC84-52 
	% -----
		\label{DC84-52}

		And whoso receiveth not my voice is not acquainted with my voice, and is not of me.

		% \begin{framed}
		% \scriptsize
			% \begin{compactdesc}
				% \item[JSP] 
				% % \item[BC] 
				% % \item[]
			% \end{compactdesc}
		% \end{framed}

	% -----
	\subsubsection{Section 84:53} % *DC84-53 
	% -----
		\label{DC84-53}

		And by this you may know the righteous from the wicked, and that the whole world groaneth under sin and darkness even now.

		% \begin{framed}
		% \scriptsize
			% \begin{compactdesc}
				% \item[JSP] 
				% % \item[BC] 
				% % \item[]
			% \end{compactdesc}
		% \end{framed}

	% -----
	\subsubsection{Section 84:54} % *DC84-54 
	% -----
		\label{DC84-54}

		And your minds in times past have been \hyperref[DARKNESS-1828]{darkened} because of unbelief, and because you have treated lightly the things you have received---

		% \begin{framed}
		% \scriptsize
			% \begin{compactdesc}
				% \item[JSP] 
				% % \item[BC] 
				% % \item[]
			% \end{compactdesc}
		% \end{framed}

	% -----
	\subsubsection{Section 84:55} % *DC84-55 
	% -----
		\label{DC84-55}

		Which vanity and unbelief have brought the whole church under condemnation.

		% \begin{framed}
		% \scriptsize
			% \begin{compactdesc}
				% \item[JSP] 
				% % \item[BC] 
				% % \item[]
			% \end{compactdesc}
		% \end{framed}

	% -----
	\subsubsection{Section 84:56} % *DC84-56 
	% -----
		\label{DC84-56}

		And this condemnation resteth upon the children of Zion, even all.

		% \begin{framed}
		% \scriptsize
			% \begin{compactdesc}
				% \item[JSP] 
				% % \item[BC] 
				% % \item[]
			% \end{compactdesc}
		% \end{framed}

	% -----
	\subsubsection{Section 84:57} % *DC84-57 
	% -----
		\label{DC84-57}

		And they shall remain under this condemnation until they repent and remember the new covenant, even the Book of Mormon and the former commandments which I have given them, not only to say, but to do according to that which I have written---%
			\footnote{%
				% \label{}%
				Is the Lord saying here that he wrote the Book of Mormon?
				}

		\begin{framed}
		\scriptsize
			\begin{compactdesc}
				\item[JSP] and thay shall remain under this condemnation until they repent and remember the new covenant even the book of Mormon and the former commandments which I have given them, not only to say but to do according to that which I have writen
				% \item[BC] 
				% \item[]
			\end{compactdesc}
		\end{framed}

	% -----
	\subsubsection{Section 84:58} % *DC84-58 
	% -----
		\label{DC84-58}

		That they may bring forth fruit \hyperref[MEET]{meet} for their Father's kingdom; otherwise there remaineth a scourge and judgment to be poured out upon the children of Zion.

		% \begin{framed}
		% \scriptsize
			% \begin{compactdesc}
				% \item[JSP] 
				% % \item[BC] 
				% % \item[]
			% \end{compactdesc}
		% \end{framed}

	% -----
	\subsubsection{Section 84:59} % *DC84-59 
	% -----
		\label{DC84-59}

		For shall the children of the kingdom pollute my holy land?  Verily, I say unto you, Nay.


		% \begin{framed}
		% \scriptsize
			% \begin{compactdesc}
				% \item[JSP] 
				% % \item[BC] 
				% % \item[]
			% \end{compactdesc}
		% \end{framed}

	% -----
	\subsubsection{Section 84:60} % *DC84-60 
	% -----
		\label{DC84-60}

		Verily, verily, I say unto you who now hear my words, which are my voice, blessed are ye inasmuch as you receive these things;

		% \begin{framed}
		% \scriptsize
			% \begin{compactdesc}
				% \item[JSP] 
				% % \item[BC] 
				% % \item[]
			% \end{compactdesc}
		% \end{framed}

	% -----
	\subsubsection{Section 84:61} % *DC84-61 
	% -----
		\label{DC84-61}

		For I will forgive you of your sins with this commandment---that you remain steadfast in your minds in solemnity and the spirit of prayer, in bearing testimony to all the world of those things which are communicated unto you.

		% \begin{framed}
		% \scriptsize
			% \begin{compactdesc}
				% \item[JSP] 
				% % \item[BC] 
				% % \item[]
			% \end{compactdesc}
		% \end{framed}

	% -----
	\subsubsection{Section 84:62} % *DC84-62 
	% -----
		\label{DC84-62}

		Therefore, go ye into all the world; and unto whatsoever place ye cannot go ye shall send, that the testimony may go from you into all the world unto every creature.

		% \begin{framed}
		% \scriptsize
			% \begin{compactdesc}
				% \item[JSP] 
				% % \item[BC] 
				% % \item[]
			% \end{compactdesc}
		% \end{framed}

	% -----
	\subsubsection{Section 84:63} % *DC84-63 
	% -----
		\label{DC84-63}

		And as I said unto mine apostles, even so I say unto you, for you are mine apostles, even God's high priests; ye are they whom my Father hath given me; ye are my friends;

		% \begin{framed}
		% \scriptsize
			% \begin{compactdesc}
				% \item[JSP] 
				% % \item[BC] 
				% % \item[]
			% \end{compactdesc}
		% \end{framed}

	% -----
	\subsubsection{Section 84:64} % *DC84-64 
	% -----
		\label{DC84-64}

		Therefore, as I said unto mine apostles I say unto you again, that every soul who believeth on your words, and is baptized by water for the remission of sins, shall receive the Holy Ghost.

		% \begin{framed}
		% \scriptsize
			% \begin{compactdesc}
				% \item[JSP] 
				% % \item[BC] 
				% % \item[]
			% \end{compactdesc}
		% \end{framed}

	% -----
	\subsubsection{Section 84:65} % *DC84-65 
	% -----
		\label{DC84-65}

		And these signs shall follow them that believe---

		% \begin{framed}
		% \scriptsize
			% \begin{compactdesc}
				% \item[JSP] 
				% % \item[BC] 
				% % \item[]
			% \end{compactdesc}
		% \end{framed}

	% -----
	\subsubsection{Section 84:66} % *DC84-66 
	% -----
		\label{DC84-66}

		In my name they shall do many wonderful works;

		% \begin{framed}
		% \scriptsize
			% \begin{compactdesc}
				% \item[JSP] 
				% % \item[BC] 
				% % \item[]
			% \end{compactdesc}
		% \end{framed}

	% -----
	\subsubsection{Section 84:67} % *DC84-67 
	% -----
		\label{DC84-67}

		In my name they shall cast out devils;

		% \begin{framed}
		% \scriptsize
			% \begin{compactdesc}
				% \item[JSP] 
				% % \item[BC] 
				% % \item[]
			% \end{compactdesc}
		% \end{framed}

	% -----
	\subsubsection{Section 84:68} % *DC84-68 
	% -----
		\label{DC84-68}

		In my name they shall heal the sick;

		% \begin{framed}
		% \scriptsize
			% \begin{compactdesc}
				% \item[JSP] 
				% % \item[BC] 
				% % \item[]
			% \end{compactdesc}
		% \end{framed}

	% -----
	\subsubsection{Section 84:69} % *DC84-69 
	% -----
		\label{DC84-69}

		In my name they shall open the eyes of the blind, and unstop the ears of the deaf;

		% \begin{framed}
		% \scriptsize
			% \begin{compactdesc}
				% \item[JSP] 
				% % \item[BC] 
				% % \item[]
			% \end{compactdesc}
		% \end{framed}

	% -----
	\subsubsection{Section 84:70} % *DC84-70 
	% -----
		\label{DC84-70}

		And the tongue of the dumb shall speak;

		% \begin{framed}
		% \scriptsize
			% \begin{compactdesc}
				% \item[JSP] 
				% % \item[BC] 
				% % \item[]
			% \end{compactdesc}
		% \end{framed}

	% -----
	\subsubsection{Section 84:71} % *DC84-71 
	% -----
		\label{DC84-71}

		And if any man shall administer poison unto them it shall not hurt them;

		% \begin{framed}
		% \scriptsize
			% \begin{compactdesc}
				% \item[JSP] 
				% % \item[BC] 
				% % \item[]
			% \end{compactdesc}
		% \end{framed}

	% -----
	\subsubsection{Section 84:72} % *DC84-72 
	% -----
		\label{DC84-72}

		And the poison of a serpent shall not have power to harm them.

		% \begin{framed}
		% \scriptsize
			% \begin{compactdesc}
				% \item[JSP] 
				% % \item[BC] 
				% % \item[]
			% \end{compactdesc}
		% \end{framed}

	% -----
	\subsubsection{Section 84:73} % *DC84-73 
	% -----
		\label{DC84-73}

		But a commandment I give unto them, that they shall not boast themselves of these things, neither speak them before the world; for these things are given unto you for your profit and for salvation.

		% \begin{framed}
		% \scriptsize
			% \begin{compactdesc}
				% \item[JSP] 
				% % \item[BC] 
				% % \item[]
			% \end{compactdesc}
		% \end{framed}

	% -----
	\subsubsection{Section 84:74} % *DC84-74 
	% -----
		\label{DC84-74}

		Verily, verily, I say unto you, they who believe not on your words, and are not baptized in water in my name, for the remission of their sins, that they may receive the Holy Ghost, shall be damned, and shall not come into my Father's kingdom where my Father and I am.

		% \begin{framed}
		% \scriptsize
			% \begin{compactdesc}
				% \item[JSP] 
				% % \item[BC] 
				% % \item[]
			% \end{compactdesc}
		% \end{framed}

	% -----
	\subsubsection{Section 84:75} % *DC84-75 
	% -----
		\label{DC84-75}

		And this revelation unto you, and commandment, is in force from this very hour upon all the world, and the gospel is unto all who have not received it.

		% \begin{framed}
		% \scriptsize
			% \begin{compactdesc}
				% \item[JSP] 
				% % \item[BC] 
				% % \item[]
			% \end{compactdesc}
		% \end{framed}

	% -----
	\subsubsection{Section 84:76} % *DC84-76 
	% -----
		\label{DC84-76}

		But, verily I say unto all those to whom the kingdom has been given---from you it must be preached unto them, that they shall repent of their former evil works; for they are to be upbraided for their evil hearts of unbelief, and your brethren in Zion for their rebellion against you at the time I sent you.

		% \begin{framed}
		% \scriptsize
			% \begin{compactdesc}
				% \item[JSP] 
				% % \item[BC] 
				% % \item[]
			% \end{compactdesc}
		% \end{framed}

	% -----
	\subsubsection{Section 84:77} % *DC84-77 
	% -----
		\label{DC84-77}

		And again I say unto you, my friends, for from henceforth I shall call you friends, it is expedient that I give unto you this commandment, that ye become even as my friends in days when I was with them, traveling to preach the gospel in my power;

		% \begin{framed}
		% \scriptsize
			% \begin{compactdesc}
				% \item[JSP] 
				% % \item[BC] 
				% % \item[]
			% \end{compactdesc}
		% \end{framed}

	% -----
	\subsubsection{Section 84:78} % *DC84-78 
	% -----
		\label{DC84-78}

		For I suffered them not to have purse or scrip, neither two coats.

		% \begin{framed}
		% \scriptsize
			% \begin{compactdesc}
				% \item[JSP] 
				% % \item[BC] 
				% % \item[]
			% \end{compactdesc}
		% \end{framed}

	% -----
	\subsubsection{Section 84:79} % *DC84-79 
	% -----
		\label{DC84-79}

		Behold, I send you out to prove the world, and the laborer is worthy of his hire.

		% \begin{framed}
		% \scriptsize
			% \begin{compactdesc}
				% \item[JSP] 
				% % \item[BC] 
				% % \item[]
			% \end{compactdesc}
		% \end{framed}

	% -----
	\subsubsection{Section 84:80} % *DC84-80 
	% -----
		\label{DC84-80}

		And any man that shall go and preach this gospel of the kingdom, and fail not to continue faithful in all things, shall not be weary in mind, neither \hyperref[DARKNESS-1828]{darkened}, neither in body, limb, nor joint; and a hair of his head shall not fall to the ground unnoticed.  And they shall not go hungry, neither athirst.

		% \begin{framed}
		% \scriptsize
			% \begin{compactdesc}
				% \item[JSP] 
				% % \item[BC] 
				% % \item[]
			% \end{compactdesc}
		% \end{framed}

	% -----
	\subsubsection{Section 84:81} % *DC84-81 
	% -----
		\label{DC84-81}

		Therefore, take ye no thought for the morrow, for what ye shall eat, or what ye shall drink, or wherewithal ye shall be clothed.

		% \begin{framed}
		% \scriptsize
			% \begin{compactdesc}
				% \item[JSP] 
				% % \item[BC] 
				% % \item[]
			% \end{compactdesc}
		% \end{framed}

	% -----
	\subsubsection{Section 84:82} % *DC84-82 
	% -----
		\label{DC84-82}

		For, consider the lilies of the field, how they grow, they toil not, neither do they spin; and the kingdoms of the world, in all their glory, are not arrayed like one of these.

		% \begin{framed}
		% \scriptsize
			% \begin{compactdesc}
				% \item[JSP] 
				% % \item[BC] 
				% % \item[]
			% \end{compactdesc}
		% \end{framed}

	% -----
	\subsubsection{Section 84:83} % *DC84-83 
	% -----
		\label{DC84-83}

		For your Father, who is in heaven, knoweth that you have need of all these things.

		% \begin{framed}
		% \scriptsize
			% \begin{compactdesc}
				% \item[JSP] 
				% % \item[BC] 
				% % \item[]
			% \end{compactdesc}
		% \end{framed}

	% -----
	\subsubsection{Section 84:84} % *DC84-84 
	% -----
		\label{DC84-84}

		Therefore, let the morrow take thought for the things of itself.

		% \begin{framed}
		% \scriptsize
			% \begin{compactdesc}
				% \item[JSP] 
				% % \item[BC] 
				% % \item[]
			% \end{compactdesc}
		% \end{framed}

	% -----
	\subsubsection{Section 84:85} % *DC84-85 
	% -----
		\label{DC84-85}

		Neither take ye thought beforehand what ye shall say; but treasure up in your minds continually the words of life, and it shall be given you in the very hour that portion that shall be meted unto every man.

		% \begin{framed}
		% \scriptsize
			% \begin{compactdesc}
				% \item[JSP] 
				% % \item[BC] 
				% % \item[]
			% \end{compactdesc}
		% \end{framed}

	% -----
	\subsubsection{Section 84:86} % *DC84-86 
	% -----
		\label{DC84-86}

		Therefore, let no man among you, for this commandment is unto all the faithful who are called of God in the church unto the ministry, from this hour take purse or scrip, that goeth forth to proclaim this gospel of the kingdom.

		% \begin{framed}
		% \scriptsize
			% \begin{compactdesc}
				% \item[JSP] 
				% % \item[BC] 
				% % \item[]
			% \end{compactdesc}
		% \end{framed}

	% -----
	\subsubsection{Section 84:87} % *DC84-87 
	% -----
		\label{DC84-87}

		Behold, I send you out to reprove the world of all their unrighteous deeds, and to teach them of a judgment which is to come.

		% \begin{framed}
		% \scriptsize
			% \begin{compactdesc}
				% \item[JSP] 
				% % \item[BC] 
				% % \item[]
			% \end{compactdesc}
		% \end{framed}

	% -----
	\subsubsection{Section 84:88} % *DC84-88 
	% -----
		\label{DC84-88}

		And whoso receiveth you, there I will be also, for I will go before your face.  I will be on your right hand and on your left, and my Spirit shall be in your hearts, and mine angels round about you, to bear you up.

		% \begin{framed}
		% \scriptsize
			% \begin{compactdesc}
				% \item[JSP] 
				% % \item[BC] 
				% % \item[]
			% \end{compactdesc}
		% \end{framed}

	% -----
	\subsubsection{Section 84:89} % *DC84-89 
	% -----
		\label{DC84-89}

		Whoso receiveth you receiveth me; and the same will feed you, and clothe you, and give you money.

		% \begin{framed}
		% \scriptsize
			% \begin{compactdesc}
				% \item[JSP] 
				% % \item[BC] 
				% % \item[]
			% \end{compactdesc}
		% \end{framed}

	% -----
	\subsubsection{Section 84:90} % *DC84-90 
	% -----
		\label{DC84-90}

		And he who feeds you, or clothes you, or gives you money, shall in nowise lose his reward.

		% \begin{framed}
		% \scriptsize
			% \begin{compactdesc}
				% \item[JSP] 
				% % \item[BC] 
				% % \item[]
			% \end{compactdesc}
		% \end{framed}

	% -----
	\subsubsection{Section 84:91} % *DC84-91 
	% -----
		\label{DC84-91}

		And he that doeth not these things is not my disciple; by this you may know my disciples.

		% \begin{framed}
		% \scriptsize
			% \begin{compactdesc}
				% \item[JSP] 
				% % \item[BC] 
				% % \item[]
			% \end{compactdesc}
		% \end{framed}

	% -----
	\subsubsection{Section 84:92} % *DC84-92 
	% -----
		\label{DC84-92}

		He that receiveth you not, go away from him alone by yourselves, and cleanse your feet even with water, pure water, whether in heat or in cold, and bear testimony of it unto your Father which is in heaven, and return not again unto that man.

		% \begin{framed}
		% \scriptsize
			% \begin{compactdesc}
				% \item[JSP] 
				% % \item[BC] 
				% % \item[]
			% \end{compactdesc}
		% \end{framed}

	% -----
	\subsubsection{Section 84:93} % *DC84-93 
	% -----
		\label{DC84-93}

		And in whatsoever village or city ye enter, do likewise.

		% \begin{framed}
		% \scriptsize
			% \begin{compactdesc}
				% \item[JSP] 
				% % \item[BC] 
				% % \item[]
			% \end{compactdesc}
		% \end{framed}

	% -----
	\subsubsection{Section 84:94} % *DC84-94 
	% -----
		\label{DC84-94}

		Nevertheless, search diligently and spare not; and wo unto that house, or that village or city that rejecteth you, or your words, or your testimony concerning me.

		% \begin{framed}
		% \scriptsize
			% \begin{compactdesc}
				% \item[JSP] 
				% % \item[BC] 
				% % \item[]
			% \end{compactdesc}
		% \end{framed}

	% -----
	\subsubsection{Section 84:95} % *DC84-95 
	% -----
		\label{DC84-95}

		Wo, I say again, unto that house, or that village or city that rejecteth you, or your words, or your testimony of me;

		% \begin{framed}
		% \scriptsize
			% \begin{compactdesc}
				% \item[JSP] 
				% % \item[BC] 
				% % \item[]
			% \end{compactdesc}
		% \end{framed}

	% -----
	\subsubsection{Section 84:96} % *DC84-96 
	% -----
		\label{DC84-96}

		For I, the Almighty, have laid my hands upon the nations, to scourge them for their wickedness.

		% \begin{framed}
		% \scriptsize
			% \begin{compactdesc}
				% \item[JSP] 
				% % \item[BC] 
				% % \item[]
			% \end{compactdesc}
		% \end{framed}

	% -----
	\subsubsection{Section 84:97} % *DC84-97 
	% -----
		\label{DC84-97}

		And plagues shall go forth, and they shall not be taken from the earth until I have completed my work, which shall be cut short in righteousness---

		% \begin{framed}
		% \scriptsize
			% \begin{compactdesc}
				% \item[JSP] 
				% % \item[BC] 
				% % \item[]
			% \end{compactdesc}
		% \end{framed}

	% -----
	\subsubsection{Section 84:98} % *DC84-98 
	% -----
		\label{DC84-98}

		Until all shall know me, who remain, even from the least unto the greatest, and shall be filled with the knowledge of the Lord, and shall see eye to eye,%
			\footnote{%
				% \label{}%
				For ``see eye to eye,'' see \hyperref[SEE-PHRASES]{here}, and for ``eye to eye,'' see \hyperref[EYE-PHRASES]{here}. The latter has more information.
				}
		and shall lift up their voice, and with the voice together sing this new song,%
			\footnote{%
				% \label{}%
				In the next verse begins a song which I call the ``Song of Zion.'' For more songs in the scriptures, see \hyperref[SONG-scriptures]{here}.
				}
		saying:

		% \begin{framed}
		% \scriptsize
			% \begin{compactdesc}
				% \item[JSP] 
				% % \item[BC] 
				% % \item[]
			% \end{compactdesc}
		% \end{framed}

	% -----
	\subsubsection{Section 84:99} % *DC84-99 
	% -----
		\label{DC84-99}

		The Lord hath brought again Zion; \\
		The Lord hath redeemed his people, Israel, \\
		According to the \hyperref[ELECTION]{election}%
			\footnote{%
				% \label{}%
				``Election'' pointing back to \hyperref[DC84-34]{verse 34}?
				}
		of grace, \\
		Which was brought to pass by the faith \\
		And covenant of their fathers. \\

		\begin{framed}
		\scriptsize
			\begin{compactdesc}
				\item[JSP] the lord hath brought again Zion the Lord hath redeemed his people Israel, according to the election of grace which was brought to pass by the faith and covenant of ther fathers,
				% \item[BC] 
				% \item[]
			\end{compactdesc}
		\end{framed}

	% -----
	\subsubsection{Section 84:100} % *DC84-100 
	% -----
		\label{DC84-100}

		The Lord hath redeemed his people;          \\
		And Satan is bound and \hyperref[TIME]{time} is no longer.   \\
		The Lord hath gathered all things in one.   \\
		The Lord hath brought down Zion from above. \\
		The Lord hath brought up Zion from beneath. \\

		\begin{framed}
		\scriptsize
			\begin{compactdesc}
				\item[JSP] the Lord hath redeemed his people and Satan is bound and time is no longer the Lord hath gathered all things in one the Lord hath brought down Zion from above the Lord hath brought up Zion from benieth
				% \item[BC] 
				% \item[]
			\end{compactdesc}
		\end{framed}

	% -----
	\subsubsection{Section 84:101} % *DC84-101 
	% -----
		\label{DC84-101}

		The earth hath travailed and brought forth her strength; \\ 
		And truth is established in her bowels;                  \\
		And the heavens have smiled upon her;                    \\
		And she is clothed with the glory of her God;            \\
		For he stands in the midst of his people.                \\

		\begin{framed}
		\scriptsize
			\begin{compactdesc}
				\item[JSP] the earth hath travailed and brought forth her strength and truth is established in her bowels and the heavens hath smiled upon her and she is clothed with the glory of her God for he standeth in the midst of his people,
				% \item[BC] 
				% \item[]
			\end{compactdesc}
		\end{framed}

	% -----
	\subsubsection{Section 84:102} % *DC84-102 
	% -----
		\label{DC84-102}

		Glory, and honor, and power, and might,			 \\ 
		Be ascribed to our God; for he is full of mercy, \\  
		Justice, grace and truth, and peace,             \\
		Forever and ever, Amen.                          \\

		\begin{framed}
		\scriptsize
			\begin{compactdesc}
				\item[JSP] glory and honor and power and might be ascribed to our god for he is full of mercy Justice grace and truth and peace for ever and ever Amen------
				% \item[BC] 
				% \item[]
			\end{compactdesc}
		\end{framed}

	% -----
	\subsubsection{Section 84:103} % *DC84-103 
	% -----
		\label{DC84-103}

		And again, verily, verily, I say unto you, it is expedient that every man who goes forth to proclaim mine everlasting gospel, that inasmuch as they have families, and receive money by gift, that they should send it unto them or make use of it for their benefit, as the Lord shall direct them, for thus it seemeth me good.

		% \begin{framed}
		% \scriptsize
			% \begin{compactdesc}
				% \item[JSP] 
				% % \item[BC] 
				% % \item[]
			% \end{compactdesc}
		% \end{framed}

	% -----
	\subsubsection{Section 84:104} % *DC84-104 
	% -----
		\label{DC84-104}

		And let all those who have not families, who receive money, send it up unto the bishop in Zion, or unto the bishop in Ohio, that it may be \hyperref[CONSECRATION]{consecrated} for the bringing forth of the revelations and the printing thereof, and for establishing Zion.

		% \begin{framed}
		% \scriptsize
			% \begin{compactdesc}
				% \item[JSP] 
				% % \item[BC] 
				% % \item[]
			% \end{compactdesc}
		% \end{framed}

	% -----
	\subsubsection{Section 84:105} % *DC84-105 
	% -----
		\label{DC84-105}

		And if any man shall give unto any of you a coat, or a suit, take the old and cast it unto the poor, and go on your way rejoicing.

		% \begin{framed}
		% \scriptsize
			% \begin{compactdesc}
				% \item[JSP] 
				% % \item[BC] 
				% % \item[]
			% \end{compactdesc}
		% \end{framed}

	% -----
	\subsubsection{Section 84:106} % *DC84-106 
	% -----
		\label{DC84-106}

		And if any man among you be strong in the Spirit, let him take with him him that is weak, that he may be \hyperref[EDIFICATION]{edified} in all meekness, that he may become strong also.

		% \begin{framed}
		% \scriptsize
			% \begin{compactdesc}
				% \item[JSP] 
				% % \item[BC] 
				% % \item[]
			% \end{compactdesc}
		% \end{framed}

	% -----
	\subsubsection{Section 84:107} % *DC84-107 
	% -----
		\label{DC84-107}

		Therefore, take with you those who are ordained unto the lesser priesthood, and send them before you to make appointments, and to prepare the way, and to fill appointments that you yourselves are not able to fill.

		% \begin{framed}
		% \scriptsize
			% \begin{compactdesc}
				% \item[JSP] 
				% % \item[BC] 
				% % \item[]
			% \end{compactdesc}
		% \end{framed}

	% -----
	\subsubsection{Section 84:108} % *DC84-108 
	% -----
		\label{DC84-108}

		Behold, this is the way that mine apostles, in ancient days, built up my church unto me.

		% \begin{framed}
		% \scriptsize
			% \begin{compactdesc}
				% \item[JSP] 
				% % \item[BC] 
				% % \item[]
			% \end{compactdesc}
		% \end{framed}

	% -----
	\subsubsection{Section 84:109} % *DC84-109 
	% -----
		\label{DC84-109}

		Therefore, let every man stand in his own office, and labor in his own calling; and let not the head say unto the feet it hath no need of the feet; for without the feet how shall the body be able to stand?

		% \begin{framed}
		% \scriptsize
			% \begin{compactdesc}
				% \item[JSP] 
				% % \item[BC] 
				% % \item[]
			% \end{compactdesc}
		% \end{framed}

	% -----
	\subsubsection{Section 84:110} % *DC84-110 
	% -----
		\label{DC84-110}

		Also the body hath need of every member, that all may be \hyperref[EDIFICATION]{edified} together, that the system%
			\footnote{%
				% \label{}%
				This is the only occurrence of ``system'' in all the scriptures---it does appear one other time but it's in a note to the JS-H. In this verse, it seems to be referring to the ``body'' as a collection of ``members'' who have different functions and differing levels of strength. We need everyone, in order for everyone to grow, and so that the system can be kept ``perfect.''
				From the 1828 dictionary, for ``system'':
					\begin{compactitem}
						\item[1] An assemblage of things adjusted into a regular whole; or a whole plan or scheme consisting of many parts connected in such a manner as to create a chain of mutual dependencies; or a regular union of principles or parts forming one entire thing. Thus we say, a \textit{system} of logic, a \textit{system} of philosophy, a \textit{system} of government, a \textit{system} of principles, the solar \textit{system} the Copernican \textit{system} a \textit{system} of divinity, a \textit{system} of law, a \textit{system} of morality, a \textit{system} of husbandry, a \textit{system} of botany or of chimistry.
						\item[2] Regular method or order.
						\item[3] In music, an interval compounded or supposed to be compounded of several lesser intervals, as the fifth octave, etc. the elements of which are called diastems.
					\end{compactitem}
				}
		may be kept \hyperref[PERFECTION]{perfect}.

		% \begin{framed}
		% \scriptsize
			% \begin{compactdesc}
				% \item[JSP] 
				% % \item[BC] 
				% % \item[]
			% \end{compactdesc}
		% \end{framed}

	% -----
	\subsubsection{Section 84:111} % *DC84-111 
	% -----
		\label{DC84-111}

		And behold, the high priests should travel, and also the elders, and also the lesser priests; but the deacons and teachers should be appointed to watch over the church, to be standing ministers unto the church.

		% \begin{framed}
		% \scriptsize
			% \begin{compactdesc}
				% \item[JSP] 
				% % \item[BC] 
				% % \item[]
			% \end{compactdesc}
		% \end{framed}

	% -----
	\subsubsection{Section 84:112} % *DC84-112 
	% -----
		\label{DC84-112}

		And the bishop, Newel K. Whitney, also should travel round about and among all the churches, searching after the poor to administer to their wants by humbling the rich and the proud.

		% \begin{framed}
		% \scriptsize
			% \begin{compactdesc}
				% \item[JSP] 
				% % \item[BC] 
				% % \item[]
			% \end{compactdesc}
		% \end{framed}

	% -----
	\subsubsection{Section 84:113} % *DC84-113 
	% -----
		\label{DC84-113}

		He should also employ an agent to take charge and to do his secular business as he shall direct.

		% \begin{framed}
		% \scriptsize
			% \begin{compactdesc}
				% \item[JSP] 
				% % \item[BC] 
				% % \item[]
			% \end{compactdesc}
		% \end{framed}

	% -----
	\subsubsection{Section 84:114} % *DC84-114 
	% -----
		\label{DC84-114}

		Nevertheless, let the bishop go unto the city of New York, also to the city of Albany, and also to the city of Boston, and warn the people of those cities with the sound of the gospel, with a loud voice, of the desolation and utter abolishment which await them if they do reject these things.

		% \begin{framed}
		% \scriptsize
			% \begin{compactdesc}
				% \item[JSP] 
				% % \item[BC] 
				% % \item[]
			% \end{compactdesc}
		% \end{framed}

	% -----
	\subsubsection{Section 84:115} % *DC84-115 
	% -----
		\label{DC84-115}

		For if they do reject these things the hour of their judgment is nigh, and their house shall be left unto them desolate.

		% \begin{framed}
		% \scriptsize
			% \begin{compactdesc}
				% \item[JSP] 
				% % \item[BC] 
				% % \item[]
			% \end{compactdesc}
		% \end{framed}

	% -----
	\subsubsection{Section 84:116} % *DC84-116 
	% -----
		\label{DC84-116}

		Let him trust in me and he shall not be confounded; and a hair of his head shall not fall to the ground unnoticed.

		% \begin{framed}
		% \scriptsize
			% \begin{compactdesc}
				% \item[JSP] 
				% % \item[BC] 
				% % \item[]
			% \end{compactdesc}
		% \end{framed}

	% -----
	\subsubsection{Section 84:117} % *DC84-117 
	% -----
		\label{DC84-117}

		And verily I say unto you, the rest of my servants, go ye forth as your circumstances shall permit, in your several callings, unto the great and notable cities and villages, reproving the world in righteousness of all their unrighteous and ungodly deeds, setting forth clearly and understandingly%
			\footnote{%
				% \label{}%
				This is the only use of ``understandingly'' in all the scriptures; it's used in a very unusual way here. If we said it now, we would mean, ``in an understanding [i.e., empathetic, compassionate] manner,'' whereas here, it means, ``in such a way as to make something understood.'' It reminds me of \hyperref[DC68-25]{DC 68:25}.
				} 
		the desolation of abomination in the last days.

		% \begin{framed}
		% \scriptsize
			% \begin{compactdesc}
				% \item[JSP] 
				% % \item[BC] 
				% % \item[]
			% \end{compactdesc}
		% \end{framed}

	% -----
	\subsubsection{Section 84:118} % *DC84-118 
	% -----
		\label{DC84-118}

		For, with you saith the Lord Almighty, I will rend their kingdoms; I will not only shake the earth, but the starry heavens shall tremble.

		% \begin{framed}
		% \scriptsize
			% \begin{compactdesc}
				% \item[JSP] 
				% % \item[BC] 
				% % \item[]
			% \end{compactdesc}
		% \end{framed}

	% -----
	\subsubsection{Section 84:119} % *DC84-119 
	% -----
		\label{DC84-119}

		For I, the Lord, have put forth my hand to exert the powers of heaven;%
			\footnote{%
				% \label{}%
				For ``powers of heaven,'' see \hyperref[POWER-PHRASES]{here}.
				}
		ye cannot see it now, yet a little while and ye shall see it, and know that I am, and that I will come and reign with my people.

		% \begin{framed}
		% \scriptsize
			% \begin{compactdesc}
				% \item[JSP] 
				% % \item[BC] 
				% % \item[]
			% \end{compactdesc}
		% \end{framed}

	% -----
	\subsubsection{Section 84:120} % *DC84-120 
	% -----
		\label{DC84-120}

		I am Alpha and Omega,%
			\footnote{%
				% \label{}%
				See \hyperref[REVELATIONNT-1-8]{Revelation 1:8} and notes there.
				} 
		the beginning and the end.  Amen.

		% \begin{framed}
		% \scriptsize
			% \begin{compactdesc}
				% \item[JSP] 
				% % \item[BC] 
				% % \item[]
			% \end{compactdesc}
		% \end{framed}